%@ pages 362-364
%@ status translated
% Chapter II, section 6: Categories with generators and exact inductive
% limits (pp. 362-364).
%
% Numbering follows the scan: this section prints LEMME 1 (the first lemma of
% the section), THEOREME 2 and THEOREME 3 (continuing THEOREME 1 of p. 356),
% PROPOSITION 13 (continuing PROPOSITION 12 of p. 361), and one corollary with
% no number at all --- hence the `unnumberedcorollary' environment.
%
% Gabriel writes the zero object as an italic `o'; throughout this
% translation it is rendered 0, the modern convention.

\section{Categories with Generators and Exact Inductive Limits}

--- In this section we suppose that $\Cat$ is an abelian category with
generators and exact inductive limits; the letter $U$ denotes a generator of
$\Cat$.

\begin{lemma}
Every subobject $M$ of an object $N$ has a complement $Q$ in $N$.
\end{lemma}

Let $E$ be the set of the subobjects $P$ of $N$ such that $M \cap P = o$. It
is enough to show that $E$ is an inductively ordered set; for this, let $F$ be
a totally ordered sub-set of $E$. Since there exists a set belonging to the
universe $\Univ$ and having the same cardinality as $F$, the supremum $S$ of
the elements $P$ of $F$ is defined. Moreover we have the formulas
\[
  S \cap M = \left(\sup_{P \in F} P\right) \cap M
           = \sup_{P \in F}\left(P \cap M\right) = o,
\]
This proves the lemma.

\begin{theorem}
Every abelian category with generators and exact inductive limits is a category
with injective envelopes.
\end{theorem}

For every object $M$ of $\Cat$ there is indeed a monomorphism from $M$ into an
injective object: this fact is proved in [10], Theorem~1.10.1. It remains valid
under our hypotheses. The theorem therefore follows from Proposition~10 and
Lemma~1.

\begin{proposition}
Let $I$ be an ordered set belonging to the universe $\Univ$; let us give
ourselves two increasing maps, $\alpha \mapsto M_\alpha$,
$\alpha \mapsto N_\alpha$, from $I$ into the set of the subobjects of an object
$P$; suppose that $M_\alpha$ is a subobject of $N_\alpha$ for every $\alpha$,
and that $N_\alpha$ is an essential extension of $M_\alpha$. Then
$\sup_\alpha N_\alpha$ is an essential extension of $\sup_\alpha M_\alpha$.
\end{proposition}

Let $Q$ be a non-zero subobject of $\sup_\alpha N_\alpha$. The formula
\[
  Q = Q \cap \sup_\alpha N_\alpha = \sup_\alpha\left(Q \cap N_\alpha\right)
\]
shows that $Q \cap N_\alpha$ is non-zero for at least one $\alpha$. It follows
that
\[
  Q \cap M_\alpha = \left(Q \cap N_\alpha\right) \cap M_\alpha
\]
is non-zero. The subobject $Q$ therefore meets $\sup_\alpha M_\alpha$ in a
non-zero subobject.

\begin{unnumberedcorollary}
If $J$ is a set belonging to $\Univ$, and if $(u_i)_{i \in J}$ is a family of
essential extensions, then $\Sigma_i u_i$ is an essential extension.
\end{unnumberedcorollary}

The corollary has already been proved when $J$ is finite. From there we pass
to the general case thanks to Proposition~13.

We end this section with an application of the preceding theorem, or rather of
Theorem~1.10.1 of [10]. The statement and its proof are due to students of
\textsc{Eilenberg}.

\begin{theorem}
Let $\Cat$ be an abelian category whose set of objects is an element of
$\Univ$. There exists an exact and faithful functor $F$ from $\Cat$ into the
category $\Ab$ of abelian groups.
\end{theorem}

Recall that an exact functor $F$ is said to be \emph{faithful} if $Ff$ is
non-zero for every non-zero morphism $f$. It comes to the same to say that $FX$
is non-zero if $X$ is a non-zero object.

We have seen that $\operatorname{Sex}(\Cat, \Ab)$ is an abelian category with
exact inductive limits. Moreover, the functor
\[
  Y \mapsto \sum_{X \in \mathbf{OC}} \operatorname{Hom}(X, Y)
\]
is a projective generator $U$ of $\operatorname{Sex}(\Cat, \Ab)$. If
$i\colon U \to F$ is a monomorphism from $U$ into an injective $F$,
Corollary~1 to Proposition~6 shows that $F$ is an exact functor. Moreover
$i(Y)$ is, for every object $Y$ of $\Cat$, an injective map from $UY$ into
$FY$; since $UY$ is non-zero, it is the same for $FY$. This proves that $F$ is
faithful.

Theorem~3 is often used to extend to abelian categories results which have
been proved for abelian groups; we give an example: let us consider in $\Cat$
the commutative and exact diagram ($\bigstar$).

\[
\tag{$\bigstar$}
\begin{tikzcd}
\phantom{A}
  & A \arrow[r, "u"] \arrow[d, "a"']
  & B \arrow[r, "v"] \arrow[d, "b"']
  & C \arrow[d, "c"'] \\
o \arrow[r]
  & A' \arrow[r, "u'"]
  & B' \arrow[r, "v'"]
  & C'
\end{tikzcd}
\]

The morphisms $u$ and $v$ induce morphisms
$u''\colon \operatorname{Ker}a \to \operatorname{Ker}b$ and
$v''\colon \operatorname{Ker}b \to \operatorname{Ker}c$. I say that the
sequence ($\bigstar\bigstar$) is exact:
\[
\tag{$\bigstar\bigstar$}
  \operatorname{Ker}a \xrightarrow{u''} \operatorname{Ker}b
  \xrightarrow{v''} \operatorname{Ker}c .
\]
Since $F$ is an exact and faithful functor, it is enough to verify the
exactness of the sequence ($\bigstar\bigstar\bigstar$):
\[
\tag{$\bigstar\bigstar\bigstar$}
  F\operatorname{Ker}a \xrightarrow{Fu''} F\operatorname{Ker}b
  \xrightarrow{Fv''} F\operatorname{Ker}c .
\]
Now the sequence ($\bigstar\bigstar\bigstar$) can be constructed from
($\bigstar\bigstar\bigstar\bigstar$) as ($\bigstar\bigstar$) was constructed
from ($\bigstar$); it is therefore enough to prove our assertion when
($\bigstar$) is a diagram of abelian groups, which is easy.

\[
\tag{$\bigstar\bigstar\bigstar\bigstar$}
\begin{tikzcd}
\phantom{FA}
  & FA \arrow[r, "Fu"] \arrow[d, "Fa"']
  & FB \arrow[r, "Fv"] \arrow[d, "Fb"']
  & FC \arrow[d, "Fc"'] \\
o \arrow[r]
  & FA' \arrow[r, "Fu'"]
  & FB' \arrow[r, "Fv'"]
  & FC'
\end{tikzcd}
\]